\section{Introduction}
\label{sec:intro}

Distributed fusion systems reduce sensor streams to the quantities on which a
decision is made. A pharmaceutical logger, for example, can reduce a
temperature history to a mean kinetic temperature and an excursion decision
\cite{usp2023mkt}; a maritime system can reduce many reports to
a track.

That reduction also removes the auditor's simplest check: retain the inputs and
recompute. An aggregator may alter a record, attach the wrong origin, omit an
inconvenient observation, or report a value that the declared operator did not
produce. Physical access can additionally expose application software and
keys. Classical verifiable computation begins with a powerful worker and a
weaker client that wants to avoid recomputation \cite{gennaro2010}. Here the
evidence producer may be the least capable device, although aggregation and
proof generation can also be delegated.

General frameworks let an application choose a verified property, while
concrete aggregation protocols prove fixed packages of properties.
Proof-carrying data uses an application-defined compliance predicate
\cite{chiesa2010pcd}; the KTV framework represents verifiability through a goal,
a judge, and an error bound \cite{cortier2016}; and aggregation protocols combine
authentication, commitments, and operator-specific checks
\cite{przydatek2003sia,palazzo2024mpvas,zhou2024ppvda}. They do not provide this
particular cumulative, run-level vocabulary for saying which part of a package
a deployment can verify.

We separate four atomic conditions---anchored manifest integrity, credential
attribution, policy-relative coverage, and computation conformance---and use
their prefixes as a cumulative reporting convention. The hierarchy is chosen
because a claim about the policy-required computation should not omit the
integrity, attribution, or coverage obligations on which that claim relies; it
is not a claim that the atomic conditions have a unique logical order. Physical
compromise is reported independently through a capture-resilience profile. No
combination proves that a sensor measured the physical world honestly.

We contribute G1--G4 with cumulative \Vzero--\Vfour{} and capture profiles; a
cross-field mechanism map; and an executable cold-chain model of nine complete
stacks, one partial baseline, and five proof circuits. The evaluation yields
design rules about representation, probabilistic coverage, and attribution
cost. The V prefix distinguishes this vocabulary from JDL fusion levels, which
grade what is inferred rather than what can be checked \cite{steinberg1999jdl}.

Generative AI assisted literature discovery, artifact development, checks, and
revision. Human authors verified all sources, code, analyses, and text and
retain full responsibility.\par
